\chapter{Modèles linéaires et quadratiques}
\label{workbook:chapter:04}
\section{Le modèle affine}
\label{workbook:section:04.01}
\begin{exerciseblock}{Le modèle affine}
\item Déterminer l'équation de la droite passant par $(2,5)$ et $(8,17)$, puis interpréter la pente.
\item Un taxi coûte $4{,}25\,$\$ au départ et $1{,}85\,$\$/km. Écrire le modèle, calculer le prix de $12$ km et résoudre $C(d)=30$.
\item Comparer $f(x)=3x-2$ et $g(x)=1{,}5x+7$: point d'intersection et domaine où $f>g$.
\end{exerciseblock}
\section{Le modèle quadratique}
\label{workbook:section:04.02}
\begin{exerciseblock}{Le modèle quadratique}
\item Mettre $f(x)=2x^2-8x+5$ sous forme canonique et donner sommet, axe, orientation et valeur minimale.
\item Une aire vaut $A(x)=x(12-x)$. Déterminer le domaine contextuel et la valeur maximale.
\item Construire une fonction quadratique ayant pour racines $-2$ et $5$ et passant par $(0,-20)$.
\end{exerciseblock}
\section{Zéros et formule quadratique}
\label{workbook:section:04.03}
\begin{exerciseblock}{Zéros et formule quadratique}
\item Résoudre $3x^2-7x-6=0$ et vérifier par somme et produit des racines.
\item Sans calculer les racines, déterminer le nombre de zéros réels de $2x^2+4x+5$ et l'interprétation graphique.
\item Trouver la valeur de $k$ pour que $x^2-6x+k=0$ ait une racine double.
\end{exerciseblock}
\section{Modéliser et interpréter}
\label{workbook:section:04.04}
\begin{exerciseblock}{Modéliser et interpréter}
\item La hauteur d'un objet est $h(t)=-4{,}9t^2+19{,}6t+1{,}5$. Déterminer hauteur initiale, instant du sommet et hauteur maximale.
\item Un modèle affine et un modèle quadratique donnent la même erreur moyenne sur des données. Proposer trois critères non numériques pour choisir entre eux.
\item Expliquer pourquoi une racine négative d'un modèle temporel peut être algébriquement correcte mais contextuellement inadmissible.
\end{exerciseblock}
\section{Voir la variation : différences premières et secondes}
\label{workbook:section:04.05}
\begin{exerciseblock}{Voir la variation : différences premières et secondes}
\item Pour le tableau $2,7,14,23,34$ associé à $x=0,1,2,3,4$, calculer les différences premières et secondes et identifier le type de modèle.
\item Montrer que pour $f(x)=ax^2+bx+c$, la différence seconde à pas $1$ vaut $2a$.
\item Un tableau a des rapports successifs presque constants mais des différences non constantes. Quel modèle faut-il tester en priorité?
\end{exerciseblock}
\section{Passer de données à un modèle}
\label{workbook:section:04.06}
\begin{exerciseblock}{Passer de données à un modèle}
\item À partir des points $(0,3)$ et $(5,18)$, construire le modèle affine puis prédire la valeur en $x=8$.
\item Les données $(0,1)$, $(1,4)$, $(2,9)$, $(3,16)$ suggèrent un modèle. Identifier celui-ci et expliquer à partir des différences.
\item Définir le résidu $r_i=y_i-\widehat y_i$ et expliquer ce qu'indique une alternance systématique de résidus positifs puis négatifs.
\end{exerciseblock}
\section{Construire, comparer et valider un modèle}
\label{workbook:section:04.07}
\begin{exerciseblock}{Construire, comparer et valider un modèle}
\item Deux modèles prédisent $48$ et $52$ pour une observation réelle de $50$. Calculer les résidus selon la convention $y-\widehat y$ et comparer les erreurs absolues.
\item Un modèle quadratique devient négatif pour de grandes valeurs de $t$ alors que la quantité représente une population. Que conclure sur son domaine de validité?
\item Proposer une démarche en cinq étapes pour valider un modèle à partir de données, d'unités et d'un contexte.
\end{exerciseblock}
\section*{Synthèse du chapitre}
\begin{synthesisbox}
\item Une entreprise observe les coûts $(0,120)$, $(10,270)$, $(20,420)$. Construire un modèle, interpréter ses paramètres et prévoir le coût de $35$ unités.
\item Un projectile est modélisé par $h(t)=-5t^2+20t+2$. Déterminer quand il atteint $17$ m, son maximum et quand il touche le sol.
\item Créer deux petits jeux de données : l'un clairement affine, l'autre clairement quadratique. Justifier par différences successives.
\end{synthesisbox}
